% % % % % % % % % % % % % % 
%
%     Skript zu NUMERIK I
%           WS14/15
%           von Prof. Dr. Blank
% Universität Regensburg
%
%
%	Config Datei: Definitionen etc.
%
% % % % % % % % % % % % % %


% ESSENTIAL PACKAGES
% ---------------------------------------------------------------
%% encodings
\usepackage[T1]{fontenc}      %8-bit font-encoding
\usepackage[utf8]{inputenc}   %8-bit input-encoding (after fontenc)
%% language
\usepackage[shorthands=off]{babel}  %language settings; no replacement of "a->ä"
%% font-settings
\usepackage{lmodern}	       %font: latin modern
\usepackage{microtype}	       %micro-typographic optimizing, e.g. ligatures

% EXTRA PACKAGES
% ---------------------------------------------------------------
%% pagestyle
\usepackage{scrpage2}		%KOMA pagestyles scrheadings, scrplain
\usepackage{enumerate}		% easy numeration
%% quotes
\usepackage{csquotes}		%easy quotes with \enquote{...}
%% maths
\usepackage{amsmath}		%improves maths-sections
\usepackage{mathtools}      %\DeclareParedDelimiter or
                            %\DefineMathOperator …
\usepackage{amsthm}		    %for proofs etc.
\usepackage{amssymb}		%maths symbols like \square for proofs
\usepackage{dsfont}		    %for \mathds{letter} to create e.g. symbol Q
\usepackage{stmaryrd}		%for math symbols like lightning bold
\usepackage{MnSymbol}
%% objects
\usepackage{wrapfig,graphicx}	%inclusion of graphics with \includegraphics{name}
\usepackage{tikz,pgf}
\usetikzlibrary{arrows}
\usetikzlibrary{matrix}
\usetikzlibrary{decorations.pathmorphing}
\usetikzlibrary{angles}
\usetikzlibrary{quotes}
\usetikzlibrary{calc}
\usetikzlibrary{positioning}

\usepackage{framed}
%% index/bibliography
\usepackage{makeidx}		%index with \makeindex in head (creates *.idx file)
\makeindex
\usepackage[backend=biber,style=alphabetic]{biblatex}	%use \printbibliography
\bibliography{numerik_script.bib}
%% (nicer)tables
\usepackage{booktabs} 	    %better spacing and lines in tables
%% font
\usepackage[normalem]{ulem}         % strike text out
%% hyperlinks
\usepackage{hyperref}
% options: colorlink=true (-> no boxes but colored hyperlinks), color of all link
\usepackage{tabularx}        %table with extendable X-column


% DEFINITIONS
% ---------------------------------------------------------------
% math (theorems)
\theoremstyle{definition}
\newtheorem{Def}{Definition}[section]
\newtheorem{Bsp}[Def]{Beispiel}
% \theoremstyle{remark}
\newtheorem{Bem}[Def]{Bemerkung}
\newtheorem{Wdh}[Def]{Wiederholung}
\theoremstyle{plain}
\newtheorem{Satz}[Def]{Satz}
\newtheorem*{satz}{Satz}
\newtheorem{Lem}[Def]{Lemma}
\newtheorem{Kor}[Def]{Korollar}
\newtheorem{Fol}[Def]{Folgerung}

% math operators
\DeclareMathOperator{\D}{D}
\DeclareMathOperator{\dd}{d}
\DeclareMathOperator{\Rang}{rang}
%\DeclareMathOperator{\Dim}{dim}
\DeclareMathOperator{\Sign}{sign}
\DeclareMathOperator{\Eps}{eps}
\DeclareMathOperator{\Tol}{tol}
\DeclareMathOperator{\Cond}{cond}
\DeclareMathOperator{\Givens}{givens}
\DeclareMathOperator{\Supp}{supp}
\DeclareMathOperator{\Span}{span}
\DeclareMathOperator{\Const}{const}
\DeclareMathOperator{\Ker}{ker}

% NEW COMMANDS
% ---------------------------------------------------------------
% symbols
\newcommand{\C}{\mathds{C}}
\newcommand{\R}{\mathds{R}}
\newcommand{\Ren}{\mathds{R}^{n}}
\newcommand{\Renn}{\mathds{R}^{n\times n}}
\newcommand{\Renm}{\mathds{R}^{n\times m}}
\newcommand{\Q}{\mathds{Q}}
\newcommand{\N}{\mathds{N}}
\newcommand{\Z}{\mathds{Z}}
\newcommand{\F}{\mathds{F}}
\newcommand{\p}{\mathcal{P}}
\newcommand{\K}{\mathcal{K}}
\newcommand{\nn}[1][\,\bullet\,]{\left\| #1 \right\|}	% norm
\newcommand{\scp}[2]{\langle #1, #2 \rangle}	% scalar prod. 
\newcommand{\ipol}[3]{#1\left( #2 \mid #3 \right)}


% for the three boxes to show float numbers
\newcommand{\floatbox}[3]{ %
  \begin{array}{|c|c|c|}
    \cline{1-3} 	
    #1 & #2 & #3\\
    \cline{1-3}
  \end{array}
}

% inserted sections a), b)
\newcommand{\extrasection}[2]{\vspace{1.5eM}\minisec{\Large\itshape \thesection #1 #2}\vspace{1eM}}

% EVIL HACKS to go along with numeration
% equation numeration
\newcommand{\sectione}[1]{ \setcounter{equation}{0}\section{#1}}
% theorem numeration
\newcommand{\subsectione}[1]{\addtocounter{Def}{1}\subsection{#1}}
\newenvironment{Satze}[1][]{ %
  \begin{Satz}[#1]  }
  { \end{Satz}
  \addtocounter{subsection}{1}}
\newenvironment{Leme}[1][]{ %
  \begin{Lem}[#1] }
  {\end{Lem}
  \addtocounter{subsection}{1}}
\newenvironment{Kore}[1][]{ %
  \begin{Kor}[#1]}
  {\end{Kor}
  \addtocounter{subsection}{1}}
\newenvironment{Fole}[1][]{ %
  \begin{Fol}[#1]}
  {\end{Fol}
  \addtocounter{subsection}{1}}
\newenvironment{Bspe}[1][]{ %
  \begin{Bsp}[#1]}
  {\end{Bsp}
  \addtocounter{subsection}{1}}
\newenvironment{Beme}[1][]{ %
  \begin{Bem}[#1]}
  {\end{Bem}
  \addtocounter{subsection}{1}}
\newenvironment{Defe}[1][]{ %
  \begin{Def}[#1]}
  {\end{Def}
  \addtocounter{subsection}{1}}
\newenvironment{Wdhe}[1][]{ %
  \begin{Wdh}[#1]}
  {\end{Wdh}
  \addtocounter{subsection}{1}}

% Pseudocode
\newenvironment{pseudocode}[1]{ %
  \begin{minipage}{#1}
    \begin{framed}
      \hspace*{1em}	
      \begin{minipage}{#1}
        \begin{tabbing}
          for~~\= for~~\= for~~\= for~~\= some really really large amount of text \= \kill
	}
	{ %
        \end{tabbing}
      \end{minipage}
      \hspace*{1em}
    \end{framed}
  \end{minipage}
}

% IMAGES
\newcounter{imagecnt}[section]

% existing images
\newenvironment{image}[1]{%

\begin{minipage}[t]{\linewidth}
\centering 
\textbf{Abb. \thesection.\arabic{imagecnt}}
\textit{~~#1}
\\[0.5eM]
}%
{
  \end{minipage}
  \stepcounter{imagecnt}\\[1eM]
}

% missing images
\newcommand{\imagemissing}[1]{
  \begin{image}{#1}
    \includegraphics[width=2cm]{images/image_missing.jpg}
  \end{image}
}	


% STYLE SETTINGS
% ---------------------------------------------------------------
% \cfoot[<text for scrplain>]{<text for scrheadings>}	
% \allowdisplaybreaks              %allow multipage for equations
\renewcommand{\theequation}{\thesection.\arabic{equation}}


%%% Local Variables:
%%% mode: latex
%%% TeX-master: "numerik_script"
%%% End:
